女性的性高潮并非只有一种模式，了解不同的高潮类型有助于去除"必须通过性交才能高潮"的焦虑。

\begin{itemize}
	\item \textbf{阴蒂高潮}：由阴蒂及其周围组织受刺激引发，是最常见、最容易达到的高潮类型。阴蒂富含神经末梢（约 8000--10000 个感觉神经末梢），且其内部结构（阴蒂脚、前庭球）可延伸至阴道前壁，因此"阴蒂"与"阴道"高潮在解剖上并非截然分开。
	\item \textbf{阴道高潮}：由阴道（尤其前壁）受刺激引发，多需较深、持续的压力。历史上曾被与"成熟程度"挂钩，现代研究认为两类高潮更多是刺激方式与个体差异的结果。
	\item \textbf{G 点高潮}：G 点（Gräfenberg spot）指阴道前壁距入口约 3--5 cm 处一片可能对刺激敏感的区域，被认为是阴蒂内部结构与尿道旁腺（Skene 腺）的复合体。并非所有女性都能找到或对之有反应，"G 点是否存在"在学界仍有争议。
	\item \textbf{混合型高潮}：多数女性通过"阴蒂 + 阴道"的复合刺激更容易达到高潮，这类刺激方式也更容易带来\textbf{连续高潮（multiple orgasms）}。
	\item \textbf{连续高潮与不应期}：部分女性能在一次活动中经历连续数次高潮，且\textbf{女性通常没有像男性那样明显的强制不应期}，或不应期极短，可在高潮后继续被刺激而再次达到高潮；也有不少女性在高潮后阴蒂高度敏感、需要短暂休息，这同样正常。
	\item \textbf{需要澄清的迷思}：不存在"只有阴道高潮才算真高潮"的标准；"同步高潮""同时达到"并非性生活的必要目标；能否高潮、高潮类型，受刺激方式、放松程度、关系与情境影响，而非"能力"或"爱得够不够"的问题。
\end{itemize}

\noindent\textit{【编者注】}关于"连续高潮"的现代理解：女性不应对期的说法有一定生理基础（女性高潮后催乳素升高幅度通常低于男性，恢复较慢的不应期多与阴蒂过度敏感有关），但个体差异极大，不必以"能来几次"衡量性生活质量。
